% SOURCE_STATEMENT: "Problem 6. Prove the existence of Broucke-H´ enon orbit for the planar 3-body problem. Do Broucke-H´ enon orbits exist for small + big masses? Is the Broucke-H´ enon solution a minimizer in the class (1, 0, 1)? (Andrea Venturelli)"
% Preparation and final self-audit: CONTRIBUTING.md prompts C and D; source/audit artifacts retained in the local physics_statements_2026_10_11 batch.
% !TeX program = lualatex
\documentclass[11pt]{article}
\usepackage[margin=25mm]{geometry}
\usepackage{fontspec}
\setmainfont{Latin Modern Roman}
\usepackage{amsmath,amssymb,mathtools}
\usepackage{unicode-math}
\setmathfont{Latin Modern Math}
\usepackage{xurl,hyperref}
\hypersetup{colorlinks=true,urlcolor=blue}
\setcounter{secnumdepth}{0}
\setlength{\parindent}{0pt}
\setlength{\parskip}{5pt}
\setlength{\emergencystretch}{3em}
\newcommand{\sref}[2]{\hyperlink{source:#1:#2}{[S#2]}}
\newcommand{\cataloguescope}[1]{\paragraph{Recorded target.}#1}
\newcommand{\R}{\mathbb R}
\newcommand{\C}{\mathbb C}
\newcommand{\N}{\mathbb N}
\newcommand{\Z}{\mathbb Z}
\newcommand{\Q}{\mathbb Q}
\newcommand{\F}{\mathbb F}
\newcommand{\eps}{\varepsilon}
\begin{document}
% UnsolvedMath ID 20002535; source catalogue code AIM-PHYSICS-0043
\setcounter{section}{20002534}
\section{Coercivity and collision boundaries for the Broucke-Henon class (1,0,1)}\label{20002535}
{\small UnsolvedMath ID 20002535 · Mathematical Physics}
\subsection{Definitions and mathematical statement}

\paragraph{Shared notation.}$\mathbb N=\{1,2,\ldots\}$; the standard number fields have their usual meanings. All additional conventions are stated locally.


\paragraph{Definitions and notation.}In the Newtonian $n$-body problem, positions $q_i\in\mathbb R^d$ and masses $m_i>0$ obey $m_i\ddot q_i=\nabla_{q_i}U$, where $U(q)=\sum_{i<j}m_im_j/|q_i-q_j|$ and the gravitational constant is normalized to one. Here $T=\tfrac12\sum_i m_i|\dot q_i|^2$, $E=T-U$, and a collision has $q_i=q_j$. The center of mass can be fixed at zero; moment of inertia is $I=\sum_i m_i|q_i|^2$. For a fixed period $P$, the action is $\mathcal A(q)=\int_0^P(T+U)\,dt$. A collision-free minimizer in an admissible loop class satisfies the Newtonian equations when variations are permitted. The period, constraints and closure of the class must be specified: a minimizing sequence can escape or approach collisions. For a planar collision-free periodic loop, identify the plane with $\mathbb C$ and define $k_{ij}=\deg(t\mapsto(q_i(t)-q_j(t))/|q_i(t)-q_j(t)|)\in\mathbb Z$. These pairwise winding numbers describe homology information; they do not determine the full braid homotopy class. The Broucke--Hénon orbit is the named planar periodic family shown in the source. Its class $(1,0,1)$ is the ordered triple of pairwise windings used there. Small and big masses concern mass ratios, not scaling all masses together.

\paragraph{Catalogue wording retained for traceability.}
Prove existence of the Broucke--Hénon orbit for the planar three-body problem. Determine whether such orbits exist for the source’s small-plus-big mass regimes. Determine whether the Broucke--Hénon solution minimizes action in class $(1,0,1)$.
\paragraph{Scope and source qualification.}All three questions remain separate. The source gives numerical and symmetry observations, including failure of coercivity from D2 symmetry alone; these are not an existence or global-minimization proof. The precise unequal-mass parameter range is not supplied in the short question. The catalogue title/status was paired with a machine-generated research attempt; that attempt is not peer-reviewed evidence of settlement. The statement covers the original source questions rather than only the title’s subcase.
\subsection{Short English statement}
Prove Broucke–Hénon periodic orbits, cover unequal small/large mass regimes, and determine minimization in class (1,0,1).
\subsection{Sources}
\begin{itemize}
\item[\hypertarget{source:20002535:1}{S1}] ULAM.AI and the UnsolvedMath Contributors, UnsolvedMath, record 20002535 (AIM-PHYSICS-0043). Catalogue curation under CC BY 4.0; definitions and English exposition added locally.
\url{https://huggingface.co/datasets/ulamai/UnsolvedMath}.
\item[\hypertarget{source:20002535:2}{S2}] American Institute of Mathematics, Variational Methods in Celestial Mechanics; notes by Kuo-Chang Chen, version 22 June 2003, source item 6.
\url{https://aimath.org/WWN/varcelest/varcelest.pdf}.
\end{itemize}
\label{end:20002535}
\end{document}
